% !TeX program = pdflatex
% =============================================================================
% Vorlage: Unterrichtsfolien (Beamer)
% Kantonsschule KFR – Physik
%
% Verwendung:
%   1. Diese Datei nach folien/<klasse>/<semester>/<thema>/ kopieren und
%      sinnvoll umbenennen, z. B.
%      folien/5-klasse/1-hs/thermodynamik/waermeleitung-folien.tex
%   2. Titel/Thema/Datum in \title/\subtitle/\date anpassen.
%   3. Folien mit \begin{frame}{Folientitel} ... \end{frame} ergänzen.
%
% Theme: bewusst kein Warsaw/Berlin/Madrid & Co. — stattdessen ein
% schlankes, selbst konfiguriertes Theme mit grosser Schrift, hohem
% Kontrast und minimalem Chrome (keine Navigationssymbole, keine
% Sidebar), damit die Folien auch von hinten im Klassenzimmer lesbar
% sind.
% =============================================================================
\documentclass[14pt,aspectratio=169]{beamer}

\usepackage[T1]{fontenc}
\usepackage[utf8]{inputenc}
\usepackage[ngerman]{babel}
\usepackage{amsmath}
\usepackage{siunitx}
\usepackage{tikz}
\usepackage{physics}
\usepackage{circuitikz}

% --- KFR-Theme: minimal, hoher Kontrast, grosse Schrift ---------------------
\usetheme{default}
\usefonttheme{professionalfonts}
\setbeamertemplate{navigation symbols}{}
\setbeamertemplate{sidebar}{}
\setbeamertemplate{headline}{}
\beamertemplatenavigationsymbolsempty

\definecolor{kfrblau}{RGB}{20,58,94}
\definecolor{kfrgrau}{RGB}{90,90,90}

\setbeamercolor{normal text}{fg=black,bg=white}
\setbeamercolor{structure}{fg=kfrblau}
\setbeamercolor{title}{fg=kfrblau,bg=white}
\setbeamercolor{frametitle}{fg=white,bg=kfrblau}
\setbeamercolor{footline}{fg=kfrgrau,bg=white}
\setbeamercolor{page number in head/foot}{fg=kfrgrau}
\setbeamercolor{itemize item}{fg=kfrblau}
\setbeamercolor{itemize subitem}{fg=kfrblau}

\setbeamerfont{frametitle}{size=\Large,series=\bfseries}
\setbeamerfont{title}{size=\LARGE,series=\bfseries}
\setbeamerfont{itemize item}{size=\normalsize}

\setbeamertemplate{itemize item}{\textbullet}
\setbeamertemplate{itemize subitem}{--}

% Schlichte Fusszeile: nur Seitenzahl, kein Autoren-/Konferenz-Ballast
\setbeamertemplate{footline}{%
  \hfill%
  \usebeamercolor[fg]{page number in head/foot}%
  \usebeamerfont{page number in head/foot}%
  \insertframenumber\,/\,\inserttotalframenumber\hspace{1em}\vspace{0.5em}
}

% Schlichter Frametitle-Balken (voller Breite, kein Untertitel-Chrome)
\setbeamertemplate{frametitle}{%
  \nointerlineskip
  \begin{beamercolorbox}[wd=\paperwidth,ht=1.4em,dp=0.4em,leftskip=0.5em]{frametitle}
    \usebeamerfont{frametitle}\insertframetitle
  \end{beamercolorbox}
}

% Werte für Titelfolie zentral setzen
\title{Newtonsche Gesetze}
\subtitle{Mechanik}
\author{}
\institute{Kantonsschule KFR}
\date{XX.XX.20XX}

\begin{document}

% --- Titelfolie --------------------------------------------------------------
\begin{frame}[plain]
  \centering
  {\usebeamerfont{title}\usebeamercolor[fg]{title}\inserttitle\par}
  \vspace{0.3em}
  {\large\insertsubtitle\par}
  \vspace{1.5em}
  {\normalsize\insertinstitute\ -- \insertdate\par}
\end{frame}

% --- Beispiel: Inhaltsfolie mit Aufzählung -----------------------------------
\begin{frame}{Die drei Newtonschen Gesetze}
  \begin{itemize}
    \item \textbf{1. Gesetz (Trägheit):} Ein Körper bleibt in Ruhe oder
      gleichförmiger Bewegung, solange keine resultierende Kraft wirkt.
    \item \textbf{2. Gesetz (Aktionsprinzip):} $F = m \cdot a$
    \item \textbf{3. Gesetz (Wechselwirkung):} Kraft $=$ Gegenkraft
      (actio $=$ reactio)
  \end{itemize}
\end{frame}

% --- Beispiel: Folie mit Platzhalter für TikZ-Diagramm/Bild ------------------
\begin{frame}{Freikörperdiagramm}
  \centering
  \begin{tikzpicture}[scale=1.2]
    % Platzhalter-Rahmen für ein Diagramm oder Bild
    \draw[kfrgrau, dashed] (-2.5,-2) rectangle (2.5,2);
    \node[kfrgrau] at (0,0) {\small Platzhalter für Diagramm/Bild};

    % Beispiel-Kraftvektor (durch echtes Diagramm ersetzen)
    \draw[-Latex, thick, kfrblau] (0,-1.8) -- (0,-0.3) node[midway, right] {$F_G$};
  \end{tikzpicture}
\end{frame}

\end{document}
